% year: תשפ"ו
% semester: א
% moed: ב
% number: 2א
% points: 12
% categories: func-analysis, derivatives
% official_solution: yes
% source: exams/מבחנים/תשפו/sol B.pdf

\begin{question}
(12 נק') חשבו תחומי קמירות ונקודות פיתול של הפונקציה
\[ f(x) = \sqrt[3]{6x + 12} \]
\end{question}

\begin{hint}
כתבו $f(x) = (6x + 12)^{1/3}$, גזרו פעמיים ובדקו את סימן $f''$. שימו לב לנקודה שבה $f''$ אינה מוגדרת.
\end{hint}

\begin{solution}
\textbf{תחום הגדרה:} שורש שלישי מוגדר לכל מספר ממשי, לכן $f$ מוגדרת ורציפה בכל $\R$. נכתוב $f(x) = (6x + 12)^{1/3}$. עבור $x \ne -2$ (כלומר $6x + 12 \ne 0$), לפי כלל השרשרת:
\[ f'(x) = \frac13 (6x + 12)^{-2/3} \cdot 6 = 2(6x + 12)^{-2/3}, \]
\[ f''(x) = 2\cdot\left(-\frac23\right)(6x + 12)^{-5/3} \cdot 6 = -8(6x + 12)^{-5/3} = \frac{-8}{(6x + 12)^{5/3}}. \]
ב-$x = -2$ הפונקציה אינה גזירה: $\frac{f(x) - f(-2)}{x + 2} = \frac{(6(x+2))^{1/3}}{x + 2} = \frac{6^{1/3}}{(x + 2)^{2/3}} \to +\infty$ (משיק אנכי).

\textbf{סימן $f''$:} החזקה $(6x + 12)^{5/3}$ היא בעלת אותו סימן כמו $6x + 12$ (חזקה אי-זוגית של שורש שלישי). לכן
\begin{itemize}
  \item $x < -2 \iff 6x + 12 < 0 \iff f''(x) > 0$, ולכן $f$ קמורה ב-$(-\infty, -2]$;
  \item $x > -2 \iff 6x + 12 > 0 \iff f''(x) < 0$, ולכן $f$ קעורה ב-$[-2, \infty)$.
\end{itemize}

\textbf{נקודת פיתול:} ב-$x = -2$ הפונקציה רציפה, יש לגרף משיק (אנכי) והקמירות מתחלפת משני צדי הנקודה. לכן $(-2, f(-2)) = (-2, 0)$ היא נקודת פיתול, אף על פי ש-$f$ אינה גזירה שם. אין נקודות פיתול נוספות, כי $f'' \ne 0$ בכל נקודה אחרת.

\textbf{תשובה:} $f$ קמורה ב-$(-\infty, -2)$, קעורה ב-$(-2, \infty)$, ונקודת הפיתול היא $(-2, 0)$.
\end{solution}
